El analisis identifico perfiles laborales diferenciados y demostro que \textbf{RF\_50\_arboles\_d7} generaliza mejor a 2026T1: MAE Q1,141.51, RMSE Q2,050.55 y $R^2$=0.4887. Su RMSE es 28.59\% menor que el de la referencia y tambien supera a LR\_sin\_regularizacion. La ventaja es compatible con la capacidad del bosque para representar no linealidades e interacciones entre educacion, ocupacion y territorio.

El desempeño no es uniforme. El mayor MAE por educacion se observo en \textbf{Doctorado}, y por dominio en \textbf{Urbano Metropolitano}. En la banda P95-P100, el ganador conserva un MAE de Q5,524.16 y un error medio de Q5,341.04, evidencia de que la cola salarial sigue siendo dificil y tiende a subestimarse. El patron debe interpretarse junto con el tamaño de cada grupo y la regresion a la media.

Las metricas son no ponderadas y describen solamente asalariados elegibles con salario positivo registrado. No representan estimaciones oficiales de Guatemala ni efectos causales. La ausencia de variables como rama de actividad, ocupacion detallada o tamaño de empresa limita la varianza explicada; un uso posterior deberia estudiar estabilidad temporal, equidad por subgrupos y estimaciones ponderadas con FACTOR bajo el diseño muestral correspondiente.
